Definieren Sie klassische versicherungsrelevante Kennzahlen in Jahresabschlüssen von VUs

Abschlussaufwendungen
Abschlussaufwendungen sind Aufwendungen, die unmittelbar oder mittelbar durch den Abschluss eines Versicherungsvertrages entstehen. 

Abschlusskostenquote
Die Abschlusskostenquote ist das Verhältnis der Abschlussaufwendungen zur Beitragssumme des Neugeschäfts. 

Kostenquote
Die Kostenquote entspricht dem Verhältnis der Personal- und Sachkosten zu den Bruttobeiträgen. 

Stornoquote
Die Stornoquote gibt den Prozentsatz der vor Vertragsablauf gekündigten oder beitragsfrei gestellten Verträge
von Versicherungen an.

Abwicklung
Die Abwicklung ist die Differenz aus in den Vorjahren gebildeten Schadenrückstellungen und 
den daraus im Berichtsjahr geleisteten Schadenzahlungen sowie den im Berichtsjahr neu gebildeten Schadenrückstellungen

Rohüberschuss
Mehrbetrag der Erträge über den Aufwendungen vor Dotierung der
Rückstellung für Beitragsrückerstattung und der Direktgutschriften sowie der Rücklagen und der Dividenden in der Lebensversicherung. 

Nichtversicherungstechnisches Ergebnis
Der Saldo aus Erträgen und Aufwendungen, die nicht direkt dem Versicherungsgeschäft zugerechnet werden können. 

Versicherungstechnisches Ergebnis
Das versicherungstechnische Ergebnis ist die Differenz aus Erträgen
und Aufwendungen aus dem reinen Versicherungsgeschäft. 

Annual Premium Equivalent
Summe aus laufenden Neugeschäftsbeiträgen und einem Zehntel der
Neugeschäfts-Einmalbeiträge. 

Anwartschaftsbarwertverfahren (Projected-Unit-Credit-Method)
Versicherungsmathematisches Bewertungsverfahren für Verpflichtungen aus betrieblicher Altersversorgung
(zum Stichtag nur Bewertung der Verpflichtungen die verdient wurden durch Arbeitnehmer)

Bedeckungsquote
Die Bedeckungsquote gibt Auskunft über das Verhältnis
zwischen den anrechnungsfähigen Eigenmitteln und
der zur Abdeckung der Risiken erforderlichen Solvenzkapitalanforderung.


Definieren Sie klassische unternehmerische Kennzahlen in Jahresabschlüssen von VUs

Gewinnzerlegung
In der Gewinnzerlegung wird der Rohüberschuss nach seinen Quellen
aufgeteilt. Somit gibt die Gewinnzerlegung im Rahmen der Nachkalkulation Auskunft darüber, woher der Überschuss stammt. Dabei wird für
jede Ergebnisquelle der tatsächliche Geschäftsverlauf den bei der Beitragsfestsetzung zugrunde gelegten Rechengrößen gegenübergestellt.

Eigenmittel
Gesamtheit des freien, unbelasteten Vermögens, welches zur Bedeckung der Solvenzkapital- und Mindestkapitalanforderung dient. 

Basiseigenmittel
Die Basiseigenmittel setzen sich gemäß § 89 Abs. 3
VAG aus dem Überschuss der Vermögenswerte über
die Verbindlichkeiten und den nachrangigen Verbindlichkeiten zusammen. 

Ausgleichsrücklage
Die Ausgleichsrücklage entspricht dem Gesamtüberschuss der Vermögenswerte über die Verbindlichkeiten
unter Abzug der sonstigen Basiseigenmittelbestandteile.

Sicherungsvermögen
Der Teil der Aktiva eines Versicherungsunternehmens, der dazu dient,
im Insolvenzfall die Ansprüche der Versicherungsnehmer zu sichern. 

Net Asset Value
Englisch für Nettoinventarwert. Wert aller materiellen und immateriellen Vermögensgegenstände 
eines Unternehmens oder Investmentfonds abzüglich sämtlicher Verbindlichkeiten. 

Beiträge
Die gebuchten Beiträge stellen den Bruttoumsatz im Prämiengeschäft
dar und beinhalten die Beiträge der Kunden zu den entsprechenden
Versicherungsprodukten. Der verdiente Beitrag beinhaltet die auf das
Geschäftsjahr entfallenden Beiträge, zuzüglich der Überträge des Vorjahres und abzüglich der Überträge in Folgejahre.

Beitragsüberträge
Bei Beitragsüberträgen handelt es sich um Beiträge für einen bestimmten Zeitraum nach dem Bilanzstichtag. 
Für diese wird eine versicherungstechnische Rückstellung im Jahresabschluss gebildet. 

Nettoverzinsung
Die Nettoverzinsung ist definiert als Quotient aus sämtlichen Erträgen
der Kapitalanlagen abzüglich der Aufwendungen für Kapitalanlagen
und dem mittleren Bestand der Kapitalanlagen zum Jahresanfang und
zum Jahresende

Stille Reserven
Nicht aus der Bilanz ersichtliche Bestandteile des Eigenkapitals von Unternehmen, die sowohl durch eine Unterbewertung von Vermögen als
auch durch eine Überbewertung von Schulden entstehen können. 

Depotforderungen/-verbindlichkeiten
Hinterlegung von Sicherheiten beim Erstversicherer durch den Rückversicherer. 


Zinszusatzreserve
Gesetzlich vorgeschriebene zusätzliche Rückstellung für Lebensversicherer, die eine vorausschauende Erhöhung der Reserven im Hinblick
auf Phasen niedriger Zinserträge vorsieht. Die Höhe der Zinszusatzreserve ist von einem Referenzzinssatz abhängig. Sinkt der Referenzzinssatz unter den Rechnungszins eines Vertrags, wird eine Zinszusatzreserve aufgebaut. 
